% ============================================================================
% PERIPHERAL sections (sensitivity / ISL / stability) + mechanism Remark +
% limitation update — STAGING. All numbers verified against the locked config
% (P_base=7, P_solar=13, B_max=18000, max-Ppeak, Q_SRC spectrum, V not pinned,
% seeds 7-16). Every load-bearing sentence was sealed in the joint review.
% ============================================================================

% ---- Mechanism Remark (insert in the methodology, after Algorithm 1) ----
\begin{remark}\label{rem_reserve}
In addition to the drift-plus-penalty rule, the proposed scheduler applies an
eclipse-reserve guard: using the deterministic orbital forecast, it declines any
action that would leave an involved satellite with less stored energy than the
platform baseline load requires through the coming eclipse (scaled by a reserve
fraction), and idles instead. The guard is part of the selection rule; the
feasible action set itself is identical for all policies.
\end{remark}

% ============================ Sensitivity ============================
\subsection{Sensitivity Analysis}\label{EV_SENS}
We sweep the scenario parameters one at a time around the nominal point, each
value over 10 seeds (Fig.~\ref{fig:sens}).

\textbf{Platform baseline (the 7B feasibility boundary).} The peak constraint
makes the platform load a hard gate for the top tier: the 7B split is feasible
only while $P_{\text{base}} + 6.37 \le 15$\,W, that is, below about 8.6\,W. The
sweep confirms the boundary empirically: at $P_{\text{base}}{=}6$\,W the
scheduler runs 842 splits with zero downtime, at 7\,W it runs about 200, at
8\,W the count collapses to 87, and at 9\,W the 7B tier drops out entirely.
Beyond 9\,W the platform load alone exceeds what the panel sustains and every
policy degrades. This boundary is set by the hardware cap and the measured 7B
peak, not by a tuned parameter, and it is the empirical basis for the per-slot
peak shield.

\textbf{Battery capacity (the scope of the safety separation).} With the
battery sized close to the eclipse requirement ($15{,}960$\,J at the nominal
load), only the proposed scheduler and the two non-splitting baselines stay
within a 5\% downtime service-level agreement, while the greedy splitters
black out; this deployability partition is invariant for
$B^{\max}\!\in\![16{,}000, 18{,}000]$\,J. It does not persist when the battery
is provisioned generously: at $B^{\max}\!\ge\!19{,}000$\,J the reckless
splitters no longer black out, and at $24{,}000$\,J Random reaches a higher
min-fill than the proposed scheduler (0.367 versus 0.296) with negligible
downtime. We report this boundary explicitly; Section~\ref{EV_LIM} analyzes
its cause.

\textbf{Other axes.} The solar-panel sweep shows the same energy story from the
supply side (the sustainability floor sits at $P^{\text{solar}}\!\approx\!12$\,W
for the 7\,W platform); under overload (arrival probability up to 0.7) every
policy's min-fill falls while the proposed scheduler's downtime stays at 0.4\%
throughout. The constellation-size sweep is discussed with the limitations: at
16 satellites the proposed scheduler's min-fill is non-monotone, and an idle
decomposition locates the direct cause in the eclipse-reserve guard of
Remark~\ref{rem_reserve}, whose blocked slots grow from 423 to 1134 (of about
7960 decision slots) while score-based idling slightly decreases; the
interaction between the per-satellite reserve requirement and the constellation
phase structure is left to the concurrent-scheduling extension.

% ============================ Dynamic ISL ============================
\subsection{Dynamic ISL Topology}\label{EV_ISL}
To make the handoff constraint~\eqref{C_pipe_isl} bind, we replace the static
topology with time-varying links of controllable availability
$p_{\text{ISL}}$ (deterministic per-window connectivity, coherence window 300
slots). As availability falls from 1.0 to 0.3 the number of 7B splits falls
from 201 to 136, so the constraint genuinely binds. The loss concentrates
entirely on the 7B-dependent source: mountain's fill drops from 0.165 to 0.154
while every other source stays flat or rises slightly as standalone capacity
flows back; min-fill falls accordingly from 0.169 to 0.157, total quality is
unchanged, and downtime stays at 0.4\%. Under heterogeneous sources, ISL
connectivity thus couples directly to the fairness of the top-tier-dependent
source, rather than to throughput or safety.

% ============================ Stability ============================
\subsection{Stability Stress Tests}\label{EV_STAB}
A twenty-orbit run (114{,}000 slots) probes the guarantees directly
(Fig.~\ref{fig:stab20}). The battery settles into a periodic orbit-scale cycle,
and the energy and quality virtual queues are mean-rate stable: their growth is
sublinear and flattening, with $Q^E(T)/T \to 7.5\times10^{-4}$ and
$Q^U(T)/T \to 1.7\times10^{-3}$ (mean backlogs reaching only 86.5 and 192.3
after twenty orbits), while blackout stays below 0.1\% of satellite-slots. This
is the empirical signature of the robust stability of Theorem~\ref{thm_robust}
and the mean-rate guarantees of Theorem~\ref{thm_stable} under periodic
zero-harvest eclipse intervals.

% ============================ Limitation update ============================
% (replace the corresponding sentences in CONTRIBUTIONS_limitation_draft.tex)
\subsection{Limitations and Future Work}\label{EV_LIM}
Our empirical advantage is confined to the energy-constrained regime, and we
state its scope precisely, in two dimensions. First, battery capacity: the
deployability separation holds for batteries sized close to the eclipse
requirement ($B^{\max}\!\in\![16{,}000,18{,}000]$\,J) and disappears when the
battery is provisioned generously, where a policy that splits at random reaches
higher min-fill and total quality than our scheduler. An idle decomposition
verifies the cause: in the generous regime the eclipse-reserve guard of
Remark~\ref{rem_reserve} is essentially inactive (zero blocked slots at
$24{,}000$\,J), and the conservatism comes entirely from the energy virtual
queue, which tracks the energy rate deficit — the standard virtual-queue
construction that guarantees mean-rate sustainability and is correct in the
constrained regime. Because it tracks a rate deficit rather than the battery
state of charge, it does not distinguish a full battery from a tight one and
therefore throttles inference that a full battery could absorb. Folding the
state of charge into the drift, a battery-aware back-pressure term, is a
natural extension that trades a more involved frame-based analysis for
cross-regime adaptivity; we leave it to future work. Second, constellation
scale: at 16 satellites the same guard blocks about 2.7 times more decision
slots than at 8, producing a non-monotone min-fill; a concurrent multi-satellite
scheduler (the simulation commits one decision per slot, a conservative special
case of the concurrent formulation) is the corresponding extension. Two further
directions follow from the same analysis: a per-resource max-min penalty that
optimizes fairness per unit of energy rather than per task, and a
higher-fidelity power model that separates the platform power system from the
compute payload rather than treating the peak cap as one shared budget.

% ---- Setup addition (insert in EV_SETUP, after the shared-feasible-set sentence) ----
% The simulation uses a serial scheduling implementation that commits one
% decision per slot, corresponding to a centralized executor in which a lead
% satellite collects state and broadcasts the decision. The abstraction is
% identical for every policy, so it does not affect the fairness of the
% comparison, but it means throughput does not scale linearly with
% constellation size. The formulation itself (constraint~\eqref{C12}) admits
% concurrent execution; the serial implementation is its conservative special
% case, and concurrent multi-satellite scheduling is a natural extension.
